\section{Kết Luận và Hướng Phát Triển}

Báo cáo đã trình bày quy trình phương pháp luận và các thuật toán phục vụ bài toán xây dựng ngữ liệu song song mức câu giữa Hán Nôm và Quốc ngữ từ tư liệu in khắc đầu thế kỷ XX. Triển khai thực nghiệm trên hai tập của bộ \textit{Sách Dạy Đấng Làm Người Sửa Mình} ($1.319$ trang hiệu lực, $636$ đơn vị ánh xạ chính văn, $5.096$ cặp câu 1--1 với $146.277$ tự hình Hán Nôm) mang lại ba đúc kết đáng chú ý cho các bài toán số hóa văn bản cổ tương tự.

Trước hết, chất lượng dóng hàng ở giai đoạn sau phụ thuộc rất lớn vào khâu chuẩn hóa vật lý và không gian mã tự hình ở đầu vào; việc phát hiện sớm các lỗi quét đúp, xoay ngược $180^\circ$, hoán vị tay sách cùng bước quy đổi $24/25$ mã PUA sang chuẩn Unicode và khử $19$ lỗi nhầm nét có hệ thống giúp loại bỏ những nguồn nhiễu gốc mà các thuật toán so khớp chuỗi đơn thuần khó tự xử lý. Thứ hai, khi nhận dạng văn bản Quốc ngữ cổ, việc dùng chính văn bản Hán Nôm đối ứng làm tập ràng buộc ngữ âm cho bước hợp nhất ROVER giữa các giả thuyết VLM và OCR truyền thống cho thấy hiệu quả rõ rệt trong việc sửa lỗi dấu phụ và nét chữ \textit{d/đ} mà vẫn giữ lại lớp chính tả lịch sử thay vì hiện đại hóa văn bản. Thứ ba, kiến trúc quy hoạch động phân cấp hai tầng với bước chiếu ranh giới câu từ Quốc ngữ sang dòng Hán Nôm liên tục ở cấp độ cụm trang đã giúp hạn chế tình trạng trượt lệch tích lũy trên văn bản thiếu dấu câu chuẩn, duy trì tỷ lệ hàng không so sánh ở mức $0{,}00\%$ và đưa độ tương đồng chính thức đạt $98{,}14\%$ trên Tập 1 cùng $98{,}45\%$ trên Tập 2.

Trong thời gian tới, quy trình này có thể tiếp tục được mở rộng theo hai hướng chính. Ở vế Hán Nôm, cơ chế đối chứng ngữ âm có thể được phát triển theo chiều ngược lại, kết hợp khoảng cách cây cấu trúc biểu ý (IDS tree-edit distance) với âm đọc của vế Quốc ngữ nhằm gợi ý xếp hạng tự động cho các tự hình bị mờ nét cục bộ. Song song đó, việc xây dựng mô hình ngôn ngữ âm tiết chuyên biệt cho tiếng Việt thế kỷ XVII--XX sẽ giúp cải thiện thêm độ mượt cú pháp và độ chính xác của bước khôi phục phiên âm trên những trang sách bị thất lạc hoặc hư hại nặng.
